% ===========================================================================
\section{ECDSA on $\secp$ and the target relation}\label{sec:overview}
% ===========================================================================

We briefly recall the ECDSA signature verification procedure on the
$\secp$ curve, then state the relation our UAIR$^+$ arithmetizes.

\paragraph{The $\secp$ curve.}
Let $p$ be the $256$-bit $\secp$ base prime and let $E/\F_p$ denote
the short Weierstrass curve $y^2 = x^3 + 7$. Let $G \in E(\F_p)$ be
the standard base point and let $n$ be its prime order (also $256$-bit,
with $p < 2n$). The cyclic subgroup $\langle G \rangle$ is the order-$n$
subgroup in which all signatures live.

\paragraph{ECDSA verification.}
A signature on a message $e \in \F_n$ under public key $Q \in \langle G \rangle$
is a pair $(r, s) \in \F_n \times \F_n^\ast$. The verifier computes
$$
u_1 := e \cdot s^{-1} \bmod n, \qquad u_2 := r \cdot s^{-1} \bmod n,
$$
forms the point
$$
R := u_1 G \;+\; u_2 Q \;\in\; E(\F_p),
$$
and accepts iff $R \neq \OOO$ and the affine $x$-coordinate $R_x$
satisfies $\mathrm{int}(R_x) \equiv r \pmod n$.

\paragraph{Jacobian coordinates.}
Throughout the arithmetization we work in Jacobian coordinates: a triple
$\jac{P}$ with $P_Z \neq 0$ represents the affine point
$(P_X / P_Z^2,\, P_Y / P_Z^3)$, and any triple with $P_Z = 0$
represents $\OOO$. Jacobian doubling and mixed addition reduce to a
small number of $\F_p$-multiplications and squarings without field
inversions, deferring the single inversion to a final affine readout.

\paragraph{Target relation.}
We arithmetize the multi-scalar elliptic-curve multiplication
underlying ECDSA verification. The remaining checks
($u_1 \cdot s = e$ and $u_2 \cdot s = r$ in $\F_n$, and the final
$R_x \equiv r \pmod n$ test) are negligible in cost and are performed
by the verifier off-protocol against the witness final-row state.
Concretely, the relation is
\begin{equation}\label{eq:rel-ecdsa}
\mathrm{REL}_{\textsc{ecdsa-msm}} \;:=\;
\left\{(\mathbf{i},\mathbf{x};\mathbf{w}) \;\middle|\;
\begin{aligned}
&\mathbf{i} = (G), \\
&\mathbf{x} = (Q,\;G+Q,\;b_1,\;b_2,\;R_{\mathrm{init}}), \\
&\mathbf{w} = (R), \\
&Q \in \langle G \rangle,\; b_1, b_2 \in \{0,1\}^{256}, \\
&R_{\mathrm{init}} \in E(\F_p) \setminus \{\OOO\}, \\
&R \;=\; R_{\mathrm{init}}^{(2^{256})} \;+\; \sum_{t=1}^{256} 2^{256-t}\bigl(b_1[t] G + b_2[t] Q\bigr).
\end{aligned}\right\}.
\end{equation}
Here $b_1, b_2$ are the bit-decompositions of $u_1, u_2$ supplied as
public input by the verifier (after the verifier computes
$u_1, u_2$ from $(e, r, s)$ off-protocol), and the
$R_{\mathrm{init}}$ is a non-identity Jacobian seed point that
side-steps the identity-aware initial step problem discussed below.

\begin{remark}[$u_1, u_2$ and $R_{\mathrm{init}}$ supplied by the verifier]
\label{rem:public-u}
We place $u_1, u_2$ — and hence their bit-decompositions $b_1, b_2$ — in
the public input rather than the witness. The verifier is responsible
for checking $u_1 \cdot s \equiv e \pmod n$, $u_2 \cdot s \equiv r \pmod n$,
and for verifying the bit decompositions
$\sum_{t=1}^{256} b_i[t] \cdot 2^{256-t} = u_i$ in $\F_n$, all
off-protocol. This shifts the two $\F_n$ multiplications and the two
bit-decomposition checks out of the arithmetization, removing the
need for an $\F_n$-trace and any $\bitpoly{32}$ columns. The cost
saving is significant; the privacy cost is that $u_1, u_2$ are
not hidden from the proof. In scenarios where any of $e, r, s$ must
remain private, $u_1, u_2$ would have to be moved back into the
witness and the corresponding well-formedness constraints reinstated.
The same applies to other components currently in the public input
(notably $Q$ and $G+Q$).
\end{remark}

\begin{remark}[Why a non-identity $R_{\mathrm{init}}$]\label{rem:r-init}
The honest scalar-multiplication algorithm starts from the Jacobian
identity $\OOO = (1 : 1 : 0)$. The mixed-addition formulas used in our
double-and-add loop (\Cref{sec:add}) become singular when their
left summand is $\OOO$ (specifically, $Z = 0$ collapses both $H = X_{\mathrm{add}} Z^2 - X$
and $R_a = Y_{\mathrm{add}} Z^3 - Y$ to zero), so the loop cannot be
initialized at $\OOO$. The implementation side-steps this by having
the verifier supply a non-identity Jacobian seed
$R_{\mathrm{init}}$ at row~$0$ and offsetting the final result by the
appropriately doubled seed, expressed in
$\mathrm{REL}_{\textsc{ecdsa-msm}}$ via the $R_{\mathrm{init}}^{(2^{256})}$
term. Replacing the seed by unified addition formulas (which handle the
identity input cleanly) is left as a follow-up.
\end{remark}

\paragraph{Off-protocol checks.}
The remaining ECDSA verification checks are performed by the verifier
off-protocol against the witness final-row Jacobian state
$R = \jac{R} \in \F_p^3$:
\begin{itemize}
\item $R \neq \OOO$, equivalently $R_Z \neq 0$;
\item $R_x := R_X \cdot R_Z^{-2} \bmod p$ (single field inversion);
\item $\mathrm{int}(R_x) \equiv r \pmod n$.
\end{itemize}
None of these are enforced in-circuit. The implementation retains the
$\col{S\_FINAL}$ selector in the column layout so that downstream
UAIRs composing with this one (e.g.\ a gluing UAIR that does enforce
the affine readout in-circuit, or one that connects the digest of a
SHA-256 UAIR to $e$) can still gate boundary constraints on the same
row.

The remainder of this document constructs a UAIR$^+$ instance
$(\mathbb{i}_{\text{uair}},\mathbb{x}_{\text{uair}};\mathbb{w}_{\text{uair}})$
that lies in $\mathrm{REL}_{\textsc{uair}^+}$ if and only if the
input witness lies in $\mathrm{REL}_{\textsc{ecdsa-msm}}$.
